% Part: first-order-logic
% Chapter: introduction
% Section: sentences

\documentclass[../../../include/open-logic-section]{subfiles}

\begin{document}

\olfileid[tr]{fol}{int}{snt}

\olsection{\usetoken{P}{sentence}}

Artık !!{structure}s ve !!{formula}s için sağlama (``içinde doğru
olma'') tanımının bir taslağına sahibiz. Fakat bu tanım ek bir öğeye,
!!a{variable} atamasına ihtiyaç duyar; bizim istediğimizse
!!{sentence}s tanımıydı. Atamalardan nasıl kurtulacağız ve
!!{sentence}s nedir?

Biçimsel mantıkla ilk tanışmanızdan !!{variable}s ile niceleyiciler
arasındaki ilişkiye dair tartışmayı muhtemelen hatırlıyorsunuz. Bir
niceleyiciyi her zaman !!a{variable} izler ve niceleyicinin uygulandığı
!!{sentence} bölümü (onun ``kapsamı'') içinde bu !!{variable}
ifadesinin niceleyici tarafından ``bağlandığını'' kabul ederiz.
!!{formula}s içinde her !!{variable} ifadesinin karşılık gelen bir
niceleyicisi olması istenmemişti; karşılık gelen niceleyicisi olmayan
!!{variable}s ``serbest'' ya da ``bağlanmamış''tır. Serbest
!!{variable}s içermeyen bütün !!{formula}s ifadelerini !!{sentence}s
olarak alacağız.

!!a{formula}~$!A$ içindeki !!a{variable} geçişinin ne zaman bağlı
olduğu, hangi niceleyicinin onu bağladığı ve ne zaman serbest olduğu
sezgisel olarak kolayca anlaşılır. Belki bunu sınamak için parantezleri
saymayı içeren bir yöntem öğrendiniz. Yine de kesin bir tanımda ısrar
etmemiz gerekir. !!{formula}s ifadelerini tümevarımla tanımladığımızdan,
!!a{formula}~$!A$ içindeki !!a{variable}~$x$ ifadesinin serbest ve
bağlı geçişlerini de tümevarımla tanımlayabiliriz. Basitleştirilmiş
dilimiz için tanım örneğin şöyle olabilir:
\begin{enumerate}
  \item $!A$ atomikse içindeki bütün $x$ geçişleri serbesttir (yani
  $\Atom{\Obj P}{x}$ içindeki $x$ geçişi serbesttir).
  \item $!A$, $\lnot !B$ biçimindeyse $\lnot !B$ içindeki bir~$x$
  geçişi ancak ve ancak~$!B$ içindeki karşılık gelen~$x$ geçişi
  serbestse serbesttir (yani~$!B$ içindeki değişkenlerin serbest
  geçişleri, $\lnot !B$ içindeki karşılık gelen geçişlerin tam
  kendisidir).
  \item $!A$, $(!B\land !C)$ biçimindeyse $(!B\land !C)$ içindeki
  bir~$x$ geçişi ancak ve ancak~$!B$ ya da~$!C$ içindeki karşılık gelen
  $x$ geçişi serbestse serbesttir.
  \item $!A$, $\lexists[x][!B]$ biçimindeyse~$!A$ içindeki hiçbir
  $x$ geçişi serbest değildir. $!A$,~$x$'ten farklı bir
  !!{variable}~$y$ için $\lexists[y][!B]$ biçimindeyse
  $\lexists[y][!B]$ içindeki bir~$x$ geçişi ancak ve ancak~$!B$
  içindeki karşılık gelen~$x$ geçişi serbestse serbesttir.
\end{enumerate}

Değişkenlerin serbest ve bağlı geçişlerini kesin olarak tanımladıktan
sonra basitçe şunu söyleyebiliriz: Serbest !!{variable}s geçişi
içermeyen her !!a{formula}, !!a{sentence} ifadesidir.

\end{document}
